% Opening material before first \section

\begin{itemize}
\item In the case of slightly supercritical branching,  the initial count of replicative individuals weighs heavily on the chance of overall process survival.

\item In the case of subcritical branching,  when conditioning on process non-death, the expected individual count tends to a constant (a quasi-stationary mean).  The value of this, $\mean{\xt| \xt>0}$, grows to infinity in the limit of the process becoming critical from below.  In the case of the continuous time birth-death process, we can describe the inverse of this conditional mean with a formula; in the case of the discrete Galton-Watson process,  we tried to find an elementary closed form for the discrete factor $A(p)$. No such elementary formula is known, and none is expected (Koenigs analysis).  
\end{itemize}





%%%%%%%%%%%%%%%%%%%%%%%%%%%%%%%%%%%%%%%%%%%%%%%%%%%%%%%%%%%%%%%%%%%%%%%%%%%%%%%%%%%%%%%%%%%%%%%%% Special 1 criticallity thing
